\documentclass[11pt,a4paper]{article}
\usepackage[utf8]{inputenc}
\usepackage[vietnamese]{babel}
\usepackage{amsmath,amssymb}
\usepackage{geometry}
\geometry{a4paper, total={180mm,260mm}, left=15mm, right=15mm, top=18mm, bottom=18mm}
\usepackage{tikz}
\usetikzlibrary{shapes.geometric, arrows.meta, positioning, calc}
\usepackage{tcolorbox}
\tcbuselibrary{skins,breakable}
\usepackage{xcolor}
\usepackage{multicol}
\usepackage{array}
\usepackage{fontawesome}
\usepackage{fancyhdr}

\pagestyle{fancy}
\fancyhf{}
\renewcommand{\headrulewidth}{0.5pt}
\renewcommand{\footrulewidth}{0.5pt}
\fancyhead[L]{\small\bfseries\color{teal}\faBook\ SỔ TAY KIẾN THỨC TOÁN 3}
\fancyhead[R]{\small\bfseries\color{teal}Lớp Toán Thầy Lĩnh -- loptoanthaylinh.vn}
\fancyfoot[C]{\small\color{gray}-- Sổ tay kiến thức: Bài 47. Làm quen với chữ số La Mã --}

\definecolor{primary}{HTML}{005F73}
\definecolor{secondary}{HTML}{0A9396}
\definecolor{accent}{HTML}{EE9B00}
\definecolor{danger}{HTML}{AE2012}
\definecolor{boxbg}{HTML}{F4F6F8}

\begin{document}

% --- BANNER TIÊU ĐỀ ---
\begin{tcolorbox}[
  enhanced,
  colback=white,
  colframe=primary,
  arc=3mm,
  boxrule=1.2pt,
  left=4mm, right=4mm, top=2.5mm, bottom=2.5mm,
  drop shadow=black!10!white
]
  \centering
  {\Large\bfseries\color{primary}\faBookmark\ BÀI 47: LÀM QUEN VỚI CHỮ SỐ LA MÃ}\\[1.5mm]
  {\small\itshape\color{black!75}Hệ thống 3 chữ số cơ bản I, V, X; quy tắc viết số từ 1 đến 20 và thực hành xem đồng hồ số La Mã trên 1 trang A4}
\end{tcolorbox}

\vspace{-2mm}


\noindent
\begin{minipage}[t]{0.49\textwidth}\vspace{0pt}
\begin{tcolorbox}[
  enhanced,
  colback=blue!4!white,
  colframe=blue!75!black,
  arc=2mm, boxrule=0.8pt,
  left=2mm, right=2mm, top=1.8mm, bottom=1.8mm,
  title={\footnotesize\bfseries\color{white}\faKey\ 1. Ba chữ số cơ bản và quy tắc ghép}
]
{\footnotesize
\begin{itemize}\setlength{\itemsep}{1.5pt}\setlength{\leftskip}{-2.5mm}
  \item \textbf{3 chữ số La Mã gốc:}
  $$\mathbf{I} = 1 \qquad \mathbf{V} = 5 \qquad \mathbf{X} = 10$$
  \item \textbf{Quy tắc cộng (viết thêm vào bên phải):}
  \begin{itemize}\setlength{\itemsep}{0.5pt}\setlength{\leftskip}{-2.5mm}
    \item Thêm $\text{I}$ vào sau: $\text{VI} = 6$; $\text{VII} = 7$; $\text{VIII} = 8$.
    \item Thêm vào sau $\text{X}$: $\text{XI} = 11$; $\text{XV} = 15$; $\text{XX} = 20$.
  \end{itemize}
  \item \textbf{Quy tắc trừ (viết bớt ở bên trái):}
  $$\text{IV} = 5 - 1 = 4 \qquad \text{IX} = 10 - 1 = 9$$
\end{itemize}
}
\end{tcolorbox}
\end{minipage}
\hfill
\begin{minipage}[t]{0.49\textwidth}\vspace{0pt}
\begin{tcolorbox}[
  enhanced,
  colback=teal!4!white,
  colframe=teal!75!black,
  arc=2mm, boxrule=0.8pt,
  left=2mm, right=2mm, top=1.8mm, bottom=1.8mm,
  title={\footnotesize\bfseries\color{white}\faClockO\ 2. Bảng đối chiếu các số từ 1 đến 20}
]
{\footnotesize
\begin{center}
\renewcommand{\arraystretch}{1.1}
\begin{tabular}{|c|c||c|c||c|c|}
  \hline
  \textbf{Số} & \textbf{La Mã} & \textbf{Số} & \textbf{La Mã} & \textbf{Số} & \textbf{La Mã} \\
  \hline
  $1$ & $\text{I}$ & $8$ & $\text{VIII}$ & $15$ & $\text{XV}$ \\
  $2$ & $\text{II}$ & $9$ & $\text{IX}$ & $16$ & $\text{XVI}$ \\
  $3$ & $\text{III}$ & $10$ & $\text{X}$ & $17$ & $\text{XVII}$ \\
  $4$ & $\text{IV}$ & $11$ & $\text{XI}$ & $18$ & $\text{XVIII}$ \\
  $5$ & $\text{V}$ & $12$ & $\text{XII}$ & $19$ & $\text{XIX}$ \\
  $6$ & $\text{VI}$ & $13$ & $\text{XIII}$ & $20$ & $\text{XX}$ \\
  $7$ & $\text{VII}$ & $14$ & $\text{XIV}$ & $21$ & $\text{XXI}$ \\
  \hline
\end{tabular}
\end{center}
}
\end{tcolorbox}
\end{minipage}

\vspace{-2.5mm}
% --- KHỐI 3: SƠ ĐỒ TƯ DUY TIKZ ---
\noindent
\begin{tcolorbox}[
  enhanced,
  colback=orange!4!white,
  colframe=orange!80!black,
  arc=2mm, boxrule=0.8pt,
  left=2mm, right=2mm, top=1.5mm, bottom=1.5mm,
  title={\footnotesize\bfseries\color{white}\faSitemap\ SƠ ĐỒ TƯ DUY: NGUYÊN LÝ CẤU TẠO CHỮ SỐ LA MÃ}
]
\centering
\begin{tikzpicture}[
  box/.style={draw, rounded corners=2pt, align=center, font=\scriptsize\bfseries, inner sep=2.5pt, minimum height=0.75cm},
  arr/.style={-{Stealth[length=2mm]}, thick}
]
  \node[box, fill=purple!15, draw=purple!80!black, text width=2.6cm] (r1) at (0, 0) {3 KÝ TỰ GỐC\\$\mathbf{I} = 1$\\$\mathbf{V} = 5$\\$\mathbf{X} = 10$};
  \node[box, fill=teal!15, draw=teal!80!black, text width=2.8cm] (r2) at (3.8, 0) {BÊN PHẢI $\implies$ CỘNG\\$\text{VI} = 5 + 1 = 6$\\$\text{VII} = 5 + 2 = 7$\\$\text{XII} = 10 + 2 = 12$};
  \node[box, fill=blue!12, draw=blue!80!black, text width=2.8cm] (r3) at (7.6, 0) {BÊN TRÁI $\implies$ TRỪ\\$\text{IV} = 5 - 1 = 4$\\$\text{IX} = 10 - 1 = 9$\\$\text{XIV} = 10 + 4 = 14$};
  \node[box, fill=orange!15, draw=orange!80!black, text width=2.4cm] (r4) at (11.0, 0) {ỨNG DỤNG\\Mặt đồng hồ\\Số thế kỉ\\Thứ tự đề mục};

  \draw[arr, purple!80!black] (r1) -- (r2);
  \draw[arr, teal!80!black] (r2) -- (r3);
  \draw[arr, blue!80!black] (r3) -- (r4);
\end{tikzpicture}
\end{tcolorbox}

\vspace{-2.5mm}
% --- KHỐI 4: VÍ DỤ & MẸO TRÁNH BẪY ---
\noindent
\begin{minipage}[t]{0.49\textwidth}\vspace{0pt}
\begin{tcolorbox}[
  enhanced,
  colback=green!4!white,
  colframe=green!60!black,
  arc=2mm, boxrule=0.8pt,
  left=2mm, right=2mm, top=1.8mm, bottom=1.8mm,
  title={\footnotesize\bfseries\color{white}\faCheckCircle\ Ví dụ thực hành chuẩn mực}
]
{\footnotesize
\textbf{1. Đọc số La Mã:}\\
-- $\text{IX}$ đọc là: \textbf{Chín}.\\
-- $\text{XIV}$ đọc là: \textbf{Mười bốn}.\\
-- $\text{XIX}$ đọc là: \textbf{Mười chín}.\\
\textbf{2. Xem đồng hồ chữ số La Mã:}\\
Đồng hồ có kim ngắn chỉ số $\text{VIII}$, kim dài chỉ số $\text{XII}$. Lúc đó là: \textbf{8 giờ đúng}.
}
\end{tcolorbox}
\end{minipage}
\hfill
\begin{minipage}[t]{0.49\textwidth}\vspace{0pt}
\begin{tcolorbox}[
  enhanced,
  colback=red!4!white,
  colframe=red!70!black,
  arc=2mm, boxrule=0.8pt,
  left=2mm, right=2mm, top=1.8mm, bottom=1.8mm,
  title={\footnotesize\bfseries\color{white}\faExclamationTriangle\ Mẹo tránh bẫy trắc nghiệm}
]
{\footnotesize
\begin{itemize}\setlength{\itemsep}{1.5pt}\setlength{\leftskip}{-2.5mm}
  \item \textbf{Bẫy nhầm giữa IV (4) và VI (6):}\\
  $\text{I}$ đứng trước $\text{V}$ $\implies 5 - 1 = \mathbf{4}$.\\
  $\text{I}$ đứng sau $\text{V}$ $\implies 5 + 1 = \mathbf{6}$.
  \item \textbf{Quy tắc lặp:}\\
  Chữ số $\text{I}$ và $\text{X}$ không bao giờ được viết liền nhau quá $\mathbf{3}$ lần (không có $\text{IIII}$ mà phải viết là $\text{IV}$). Chữ số $\text{V}$ tuyệt đối không lặp lại!
\end{itemize}
}
\end{tcolorbox}
\end{minipage}


\end{document}
